
\nonum{section}{Lab 1: Combinational Circuits Part I}

\setcounter{section}{1}
\nonum{subsection}{Homework Assignments}
\nonum{subsubsection}{1A}
I will not be using Questasim for the labs.

\nonum{subsubsection}{1B.1}

\begin{align*}
	W &= \left(AB\right)' \\
	X &= \left(WA\right)' = \left(\left(AB\right)'A\right)' \\
	Y &= \left(WB\right)' = \left(\left(AB\right)'B\right)' \\
	Z &= \left(XY\right)' = \left(\left(\left(AB\right)'A\right)' \cdot \left(\left( AB \right)' B \right)'\right)' \\ 
	  &= AB' + A'B = A \oplus B
\end{align*}

Thus, $W$ is an XOR operation between $A$ And $B$

\nonum{subsubsection}{1B.2}

Delay time: $t_p = 10 \si{ns} = 10^{-8}\si{s}$

\begin{figure}[H]
	\begin{center}
		\begin{tikztimingtable}[
			timing/slope = 0,
			timing/rowdist = 3*1.6ex,
			timing/coldist = 2.2em,
			timing/yunit = 1.6ex,
			timing/xunit = 1.6ex,
			% x = 0.35mm,
			]
			$A$ & 2L 10H 15L 10H 2L \\
			$B$ & 7L 10H 5L 10H 7L \\
			$W$ & 8H 5L 15H 5L 6H \\
			$X$ & 3H 6L 19H 1L 5H 4L 1H \\
			$Y$ & 8H 1L 5H 4L 5H 6L 4H 6L \\
			$Z$ & 4L 6H 5L 4H 5L 6H 4L 5H \\
			\extracode
			\fulltablegrid[draw=black!70, line width=0.16ex, line cap=round,
				dash pattern=on 0pt off 1.6ex, xstep=1, ystep=1];
				\begin{scope}[semithick]
					\draw[<->] (10,-16) -- ++(1,0)
					node[midway, below, font=\scriptsize] {$10\si{ns}$};
				\end{scope}
		\end{tikztimingtable}

	\end{center}
	\caption{Timing diagram}\label{fig:timing}
\end{figure}

\nonum{subsubsection}{1C}

\begin{table}
	\begin{center}
		\begin{tabular}{| c c c c | c |}
			\hline
			$A$ & $B$ & $C$ & $D$ & $X$ \\ 
			\hline
			0 & 0 & 0 & 0 & 0 \\
			0 & 0 & 0 & 1 & 0 \\
			0 & 0 & 1 & 0 & 0 \\
			0 & 0 & 1 & 1 & 1 \\
			0 & 1 & 0 & 0 & 0 \\
			0 & 1 & 0 & 1 & 1 \\
			0 & 1 & 1 & 0 & 0 \\
			0 & 1 & 1 & 1 & 1 \\
			1 & 0 & 0 & 0 & 0 \\
			1 & 0 & 0 & 1 & 0 \\
			1 & 0 & 1 & 0 & 0 \\
			1 & 0 & 1 & 1 & 1 \\
			1 & 1 & 0 & 0 & 0 \\
			1 & 1 & 0 & 1 & 1 \\
			1 & 1 & 1 & 0 & 0 \\
			1 & 1 & 1 & 1 & 0 \\
			\hline
		\end{tabular}
		\caption{Truth Table}
		\label{tab:1b-truth-table}
	\end{center}
\end{table}

\begin{figure}[H]
	\centering
	\begin{karnaugh-map}[4][4][1][][][][]
		\minterms{5,7,12,13,14}
		\maxterms{0,1,2,3,4,5,6,7,8,9,10,11,12,13,14,15}
		% \maxterms{0,1,2,3,4,6,8,9,10,11,15}
		\implicantedge{12}{12}{14}{14}
		\implicant{5}{7}
		\implicant{5}{13}

		\draw[decorate, decoration={brace, amplitude=5pt, raise=1pt}]
			(2,5) -- (4,5) node[midway, yshift=0.35cm] {\small $A$};
		\draw[decorate, decoration={brace, amplitude=5pt, raise=1pt}]
			(3,-0.5) -- (1,-0.5) node[midway, yshift=-0.35cm] {\small $B$};
		\draw[decorate, decoration={brace, amplitude=5pt, mirror, raise=1pt}]
			(-1,2) -- (-1,0) node[midway, xshift=-0.4cm] {\small $C$};
		\draw[decorate, decoration={brace, amplitude=5pt, mirror, raise=1pt}]
			(4.5,1) -- (4.5,3) node[midway, xshift=0.4cm] {\small $D$};
	\end{karnaugh-map}
	\caption{Karnaugh map}\label{fig:kmap}
\end{figure}

This gives: $\displaystyle BC'D + A'BD + B'CD$
I do not use linux, so the Espresso GUI is not available for me.
Instead I will use a modernized CLI version~\cite{Hadipour2025espresso}:

Running in espresso :
\begin{listing}[H]
	\inputminted[breaklines, bgcolor=lightgray]{text}{resources/1c.pla}
	\caption{Espresso PLA input}\label{lst:1c-espresso}
\end{listing}

Outputs:
\begin{minted}[bgcolor=lightgray]{text}
x = (!a&c&d) | (b&!c&d) | (!b&c&d);
\end{minted}

So:
\[
	X = A'CD + BC'D + B'CD
\]

The logic espressions are slightly different, but cover the same spaces in the karnaugh map, so the expressions are logically equivalent.

\nonum{subsection}{Lab assignment 1A: Introduction SV and QuestaSim}

\nonum{subsection}{Lab assignment 1B: Spikes}

\nonum{subsection}{Lab assignment 1C: Minimization and implementation}
